\scanpage{84}
N.B. Accouplons $\mathbf Z^r$ et $\mathbf Z^r$ par
\[
((n_i),(n'_i))=\sum\delta_i n_i n'_i
\]
Alors $\lambda$ et $\mu$ sont formellement transposés par cet accouplement,
\deleted{\unsure{parfait}}, donc $\operatorname{Ker}\mu\simeq E^*_{(s)}$ est accouplé à $\operatorname{Coker}\lambda\simeq E_{(s)}$, par la formule
\[
(\alpha(x),y^*)=(x,\beta y^*)\qquad x\in\unsure{\mathbf Z^r},\ y^*\in E^*_{(s)}.
\]
\[
E_{(s)}\times E^*_{(s)}\longrightarrow\mathbf Z
\]
\deleted{Par suite,} Je dis que l'on a aussi
\[
(x,u(y))=(y,u(x))\qquad\text{pour }x,y\in E_{(s)}.
\]
En effet, on peut supposer que
\[
x=\alpha(\underbrace{(n_i)}_{\underline n}),\quad y=\alpha(\underbrace{(n'_j)}_{\underline n'}),
\]
et les deux mbs sont resp.
\[
\begin{aligned}
(\alpha(\underline n),u(y))&=(\underline n,\beta u\alpha(\underline n'))\\
(\alpha(\underline n'),u(x))&=(\underline n',\beta u\alpha(\underline n))
\end{aligned}
\]
Or on a
\[
(\underline n,\beta u\alpha(\underline n'))=\sum n_i n'_j(C_i,C_j)
\]
qui est bien symétrique en \deleted{\unsure{et}} $\underline n,\underline n'$.

Si les $\delta_i$ sont tous $1$,\inserted{(p.ex $\kappa(s)$ p-fait)} alors on trouve pour l'accouplement précédent
\[
E^*_{(s)}\simeq\operatorname{Hom}(E_{(s)},\mathbf Z)
\]
et la \unsure{relation|dualité} symétrique si on a de plus $d=1$ i.e. $d^t=1$.
Dans ce cas cette dualité, et le fait que $u:E_{(s)}\unsure{-|\to}E^*_{(s)}$ est autoadjoint, donne pour le $\operatorname{Coker}u=\Psi_{(s)}$ une \uline{autodualité} canonique (à valeurs dans $\mathbf Q/\mathbf Z$). Cette autodualité \deleted{\unsure{est liée à}}\inserted{exprime} l'obstruction à étendre la biextension qu'on a de $(Q'_\eta,Q'_\eta)=(P^0_\eta,P^0_\eta)$ par $G_{m\eta}$, par une biextension de $Q',Q'$ par $G_m$.
\margin{\figure{\begin{tikzpicture}[x=1cm,y=1cm]
\node[draw,ellipse,align=center] (a) at (1.2,0) {modèles\\de Néron};
\draw[->] (.2,1)..controls(.2,.5)and(.4,.4)..(.65,.35);
\draw (.9,1)--(1.15,.45);
\end{tikzpicture}}}
\margin{83}
